\chapter{Use the Channel Exactly as Promised}

\section*{5.1}
$\Gamma$ contains reusable assumptions; $\Delta$ contains linear obligations whose duplication would authorize two consumers of one endpoint.

\section*{5.2}
Domains of left and right are disjoint; their union equals the original domain; each binding retains its exact session type.

\section*{5.3}
$0$ is typed with the empty linear environment.

\section*{5.4}
If $\Delta=\{x:S,y:T\}$, the split $\Delta_1=\{x:S\}$ and $\Delta_2=\{x:S,y:T\}$ duplicates $x$ and is invalid.

\section*{5.5}
The split $\{x:S\}\uplus\varnothing$ drops $y:T$ from $\{x:S,y:T\}$.

\section*{5.6}
From $\{x:S,y:T,z:U\}$, one valid split is $\{x:S\}\uplus\{y:T,z:U\}$.

\section*{5.7}
Linearity prevents local misuse of endpoint obligations. It does not constrain the relative acquisition order of several otherwise well-typed sessions.

\section*{5.8}
The source theorem does not model the foreign runtime's alias creation. A separate compiler/FFI simulation theorem would be required.

\section*{5.9}
From a typing of $P\mid Q$ obtain a disjoint split $\Delta_1\uplus\Delta_2$. Swap the two subderivations and use commutativity of disjoint union.

\section*{5.10}
Induct on reduction: principal communication uses T4.3/T4.4; context cases use the induction hypothesis and rebuild splitting/restriction; congruence closure uses T4.22.

\section*{5.11}
Attach a unique runtime endpoint identity and assert that ownership-transfer operations move rather than copy that identity; reject duplicated live owners.

\section*{5.12}
A shared-channel language must change the resource discipline, for example by unrestricted/shared capabilities plus synchronization rules. The linear split cannot simply be retained unchanged.
